\documentclass[10pt,a4paper]{amsart}
\usepackage[margin=19mm]{geometry}
\usepackage{amsmath,amssymb,mathtools}
\usepackage[scr=rsfs,cal=euler]{mathalfa}
\usepackage{xfrac}
\usepackage[unicode,hidelinks,bookmarks=false]{hyperref}
\newcommand{\ie}{i.e.\ }
\newcommand{\cf}{cf.\ }
\newcommand{\resp}{respectively\ }
\newcommand{\Subsection}{\subsection}
\providecommand{\nobreakdash}{\nobreak\hbox{-}}
\setlength{\emergencystretch}{2em}
\multlinegap0pt
\usepackage{amssymb}
\usepackage{xcolor}
\usepackage{mathtools}
\usepackage[shortlabels]{enumitem}
\newtheorem{thm}{Theorem}
\newtheorem{lem}{Lemma}
\newtheorem{Def}{Definition}
\newtheorem{rem}{Remark}
\newtheorem{exmp}{Example}
\title{Mass transport on LQ optimal control problems and interpolation inequalities}
\author{Luca Rizzi, Alec Jacopo Almo Schiavoni Piazza}
\date{Source: arXiv:2605.23608v1 (CC BY 4.0)}
\begin{document}
\maketitle

\noindent\emph{Source and licence.} Mass transport on LQ optimal control problems and interpolation inequalities. Version arXiv:2605.23608v1 (CC BY 4.0), distributed by arXiv under the Creative Commons Attribution 4.0 International licence, http://creativecommons.org/licenses/by/4.0/, as recorded in the arXiv licence metadata for this exact version and shown on the abstract page https://arxiv.org/abs/2605.23608v1. Full licence notice: \texttt{licenses/arxiv-2605.23608-CC-BY-4.0.txt}.\par\smallskip

\noindent\textit{Editorial scope.} Counted unit: the article's thm dist\_displ\_interp (source0.tex lines 2471--2479). The article's complete original proof of it is delivered with it (source0.tex lines 2480--2533, one \texttt{proof} environment of 54 lines). No mathematics is written, repaired or re-derived: the statement and the proof are the article's own lines, in the article's own notation, and no retained line is substituted or re-flowed.  Retained uncounted as the source-local context the proof needs: lines 286--288; lines 1719--1727; lines 2383--2394; lines 2395--2402; lines 2435--2460.  Nothing was omitted from any retained span, so the package carries no omission record.  No retained line contains a citation, so the wrapper carries no bibliography and no bibliography entry.  No retained line contains a figure environment, an image-inclusion command or drawing code, so no image is delivered and the package declares no asset dependency.  Editorial formatting only, in addition: an AMS \texttt{amsart} preamble that restates the article's own declarations and macros used by the retained lines, namely the environments \texttt{thm}, \texttt{lem}, \texttt{Def}, \texttt{rem}, \texttt{exmp} on separate counters, together with the standard packages those lines need.  The retained environments print in this document as Theorem 1, Lemma 1, Def 1, Remark 1, Exmp 1 in order of appearance, whereas the article prints the same items with the numbering of its own full document.  The complete native source of the article as distributed in the arXiv e-print for this exact version is bundled unchanged as \texttt{sources/201225/source0.tex}.
			\begin{equation}\label{density_interpolation_inequality_formula_intro}
				\frac{1}{\rho_\tau(T_\tau(x))^{1/n}}\geq\frac{\beta_{T-\tau}^{1/n}}{\rho_0(x)^{1/n}} +\frac{\beta_{\tau}^{1/n}}{\rho_T(T(x))^{1/n}},
			\end{equation}

	\begin{thm}[Absolute continuity]\label{absolute_continuity-dyn_opt_coupl}
      Let $(A,B,Q,T)$ be a LQ optimal control problem satisfying the Kalman condition, with $T<t^*$. Let $\mu, \nu\in P_2(\mathbb{R}^n)$, with $\mu\ll\mathscr{L}^n$, and let $(\mu_\tau)=(e_{\tau\sharp}\Pi)\in\mathcal{C}([0,T],P_2(\mathbb{R}^n))$ be the unique displacement interpolation joining them. Then, $\mu_\tau\ll\mathscr{L}^n$, for every $\tau\in[0,T)$. 

        Additionally, let $\psi$ be a Kantorovich potential from $\mu$ to $\nu$ for the cost $c^{0,T}$. Then, the optimal transport map $T_\tau:\mathbb{R}^n\to\mathbb{R}^n$, from $\mu$ to $\mu_\tau$ for the cost $c^{0,\tau}$ is given by
        \begin{equation}
			T_\tau(x)=\exp_{x,\tau}(\nabla\psi(x)),\qquad\mu\text{-almost every }x\in\mathbb{R}^n, \quad\forall \tau\in[0,T].
		\end{equation}
        In particular, $T_\tau$ is $\mu$-almost everywhere differentiable.
	\end{thm}

		\begin{lem}\label{Jacobian_and_densities}
			Let $\rho\in L^1(\mathbb{R}^n)$ be a non-negative function. Let $f:\mathbb{R}^n\to\mathbb{R}^n$ be a measurable function. Suppose that there exists a Borel set $\Sigma\subset\mathbb{R}^n$ such that
			\begin{itemize}
				\item $f$ is differentiable at any point $x\in\Sigma$,
				\item $f|_{\Sigma}$ is injective,
				\item the set $\{\rho>0\}\setminus\Sigma$ is $\mathscr{L}^n$-negligible.
			\end{itemize}
			Then, the measure $f_{\sharp}(\rho\mathscr{L}^n) \ll \mathscr{L}^n$ if and only if $\det(\nabla f(x))>0$ for $\rho\mathscr{L}^n$-almost every $x\in\mathbb{R}^n$. In such case, letting $f_{\sharp}(\rho\mathscr{L}^n)=\rho_f\mathscr{L}^n$, we have
			\begin{equation}
				\rho_f(f(x))=\frac{\rho(x)}{\det(\nabla f(x))},\qquad\forall x\in\Sigma.
			\end{equation}
		\end{lem}

		\begin{thm}[Interpolation inequalities]\label{density_interpolation_inequality}
Let $(A,B,Q,T)$ be a LQ problem satisfying the Kalman condition, with $T<t^*$. Consider $\mu,\nu\in P_2(\mathbb{R}^n)$, with $\mu, \nu \ll\mathscr{L}^n$, and let $(\mu_\tau)=(e_{\tau\sharp}\Pi)$ be the unique displacement interpolation between them. Let $T_\tau$ be the optimal transport map from $\mu$ to $\mu_\tau$, for the cost $c^{0,\tau}$. Then, for $\mu$-almost every $x\in\mathbb{R}^n$ and every $\tau\in[0,T]$, we have
			\begin{equation}\label{density_interpolation_inequality_formula}
				\frac{1}{\rho_\tau(T_\tau(x))^{1/n}}\geq\frac{\beta_{T-\tau}^{1/n}}{\rho_0(x)^{1/n}} +\frac{\beta_{\tau}^{1/n}}{\rho_T(T(x))^{1/n}},
			\end{equation}
			where $\beta_s=\beta_s^{(A,B,Q,T)}$ is the distortion coefficient of the LQ problem and $\mu_\tau=\rho_\tau\mathscr{L}^n$. 
            If $\nu$ is not absolutely continuous, \eqref{density_interpolation_inequality_formula_intro} holds for $\tau\in[0,T)$ and omitting the second term in the right hand side.
            \end{thm}

		\begin{Def}[Displacement convexity classes]\label{def:DCN}
			For $N\in[1,\infty]$, the displacement convexity class of dimension $N$, denoted by $\mathcal{DC}_N$, is defined as the set of continuous convex functions $U:\mathbb{R}^{+}\to\mathbb{R}$, such that 
			\begin{itemize}
				\item $U(0)=0$,
				\item $U|_{(0,+\infty)}\in\mathcal{C}^2\big((0,+\infty)\big),$
				\item the function $\delta\mapsto u(\delta)$, defined as
			 	\begin{equation}
			 		u(\delta):=\begin{cases}
			 			\delta^NU(\delta^{-N}), &\text{if }N<\infty,\delta>0,\\
			 			e^{\delta}U(e^{-\delta}), &\text{if }N=\infty,\delta\in\mathbb{R},
			 			\end{cases}
			 	\end{equation}
                is convex.
			\end{itemize}
		\end{Def}
		\begin{rem}\label{nonincreasing_u}
			Since $U$ is convex and $U(0)=0$, the function $u$ in Definition \ref{def:DCN} is  non-increasing.
		\end{rem}
		\begin{exmp}
			Note that 
			\begin{itemize}
				\item for any $\alpha\geq 1$, the function $U(r)=r^{\alpha}$ belongs to all classes $\mathcal{DC}_N$.
				\item \label{ex_DC_2} if $\alpha<1$, then the function $U(r)=-r^{\alpha}$ belongs to $\mathcal{DC}_N$ if and only if $N\leq (1-\alpha)^{-1}$, i.e.\ $\alpha\geq 1-1/N$. In particular, the function $U(r)=-r^{1-1/N}$ is, in some sense, the minimal representative of $\mathcal{DC}_N$.
				\item the function $U_{\infty}(r)=r\log(r)$ belongs to $\mathcal{DC}_{\infty}$. It can be seen as the limit of the functions $U_N(r)=-N(r^{1-1/N}-r)$, which are the same (up to multiplication and addition of a linear function) as the functions appearing in item \eqref{ex_DC_2}. 
			\end{itemize}
		\end{exmp}

		\begin{thm}[Entropic inequalities]\label{dist_displ_interp}
	 	Let $(A,B,Q,T)$ be a LQ optimal control problem satisfying the Kalman condition, with $T<t^*$.
	 	Consider $U\in\mathcal{DC}_n$ and $\mu,\nu\in P_2(\mathbb{R}^n)$, with $\mu,\nu\ll \mathscr{L}^n$. Let $(\mu_\tau)$ be the unique displacement interpolation between $\mu$ and $\nu$. Assume that $(\mu_\tau)\subset\mathrm{Dom}(\mathcal{U})$ for every $\tau\in[0,T]$. Then, the entropy functional $\mathcal{U}$ satisfies
			\begin{multline}
				\mathcal{U}(\mu_\tau)\leq\left(\frac{T}{T-\tau}\right)^{n-1}\beta_{T-\tau}\int_{\mathbb{R}^n}U\bigg(\frac{\rho_0(x)}{\beta_{T-\tau}}\left(\frac{T-\tau}{T}\right)^n\bigg)d\mathscr{L}^n(x)\\
				+\left(\frac{T}{\tau}\right)^{n-1}\beta_{\tau}\int_{\mathbb{R}^n}U\bigg(\frac{\rho_T(y)}{\beta_{\tau}}\left(\frac{\tau}{T}\right)^n\bigg)d\mathscr{L}^n(y),\qquad\forall \tau\in(0,T).
			\end{multline}
            where $\beta_s=\beta_s^{(A,B,Q,T)}$ is the distortion coefficient of the LQ problem and $\mu_\tau=\rho_\tau\mathscr{L}^n$.
		\end{thm}

		\begin{proof}
			By Theorem \ref{absolute_continuity-dyn_opt_coupl}, $\mu_\tau\ll\mathscr{L}^n$, for every $\tau\in (0,T)$, with $\rho_\tau=d\mu_\tau/d\mathscr{L}^n$. Moreover, let $T_\tau$ be the optimal transport map between $\mu=\mu_0$ and $\mu_\tau$ for the cost $c^{0,\tau}$.
            
			Then $T_\tau$ is differentiable $\mu$-almost everywhere and by Lemma \ref{Jacobian_and_densities}, we have 
			\begin{equation}
				\det(\nabla T_\tau(x))=\frac{\rho_0(x)}{\rho_\tau(T_\tau(x))},\qquad\forall \tau\in[0,T],\quad\mu\text{-almost every }x\in\mathbb{R}^n.
			\end{equation}
			As a consequence, by change of variable formula applied to $z=T_\tau(x)$, we have 
			\begin{align}
				\int_{\mathbb{R}^n}U\big(\rho_\tau(z)\big)d\mathscr{L}^n(z)&=\int_{\mathbb{R}^n}U\big(\rho_\tau(T_\tau(x))\big)\det(\nabla T_\tau(x))d\mathscr{L}^n(x)\\
				&=\int_{\mathbb{R}^n}U\bigg(\frac{\rho_0(x)}{\det(\nabla T_\tau(x))}\bigg)\det(\nabla T_\tau(x))d\mathscr{L}^n(x).
			\end{align}
			Since $U(0)=0$, the contribution of $\{\rho_0=0\}$ is negligible, yielding 
			\begin{align}
				\mathcal{U}(\mu_\tau)&=\int_{\mathbb{R}^n}U\bigg(\frac{\rho_0(x)}{\det(\nabla T_\tau(x))}\bigg)\det(\nabla T_\tau(x))d\mathscr{L}^n(x)\\
			%	&=\int_{\{\rho_0>0\}}U\bigg(\frac{\rho_0(x)}{\det(\nabla T_t(x))}\bigg)\frac{\det(\nabla T_t(x))}{\rho_0(x)}\rho_0(x)d\mathscr{L}^n(x)\\
				&=\int_{\{\rho_0>0\}}U\bigg(\frac{\rho_0(x)}{\det(\nabla T_\tau(x))}\bigg)\frac{\det(\nabla T_\tau(x))}{\rho_0(x)}d\mu(x)\\
				&=\int_{\{\rho_0>0\}}U\big(\delta(\tau,x)^{-n}\big)\delta(\tau,x)^{n}d\mu(x)
                =\int_{\{\rho_0>0\}}u\big(\delta(\tau,x)\big)d\mu(x),
			\end{align}
			where $u(\delta)=\delta^{n}U(\delta^{-n})$ and 
			\begin{equation}
				\delta(\tau,x)=\frac{1}{\rho_\tau(T_\tau(x))^{1/n}}=\frac{\det(\nabla T_\tau(x))^{1/n}}{\rho_0(x)^{1/n}}.
			\end{equation} 
			By Theorem \ref{density_interpolation_inequality}, we have that
			\begin{equation}
				\delta(\tau,x)\geq\frac{\beta_{T-\tau}^{1/n}}{\rho_0(x)^{1/n}} +\frac{\beta_{\tau}^{1/n}}{\rho_T(T(x))^{1/n}},
			\end{equation}
			whilst, since $U\in\mathcal{DC}_n$, the function $u$ is convex and non-increasing (see Definition \ref{def:DCN} and Remark \ref{nonincreasing_u}). Therefore, continuing the computation, we have
			\begin{align}
				\mathcal{U}(\mu_\tau)&=\int_{\{\rho_0>0\}}u\left(\delta(\tau,x)\right)d\mu(x)\\
				&\leq\int_{\{\rho_0>0\}}u\left(\frac{\beta_{T-\tau}^{1/n}}{\rho_0(x)^{1/n}} +\frac{\beta_{\tau}^{1/n}}{\rho_T(T(x))^{1/n}}\right)d\mu(x)\\
				&\leq \frac{T-\tau}{T}\int_{\{\rho_0>0\}}u\bigg(\frac{T}{T-\tau}\frac{\beta_{T-\tau}^{1/n}}{\rho_0(x)^{1/n}}\bigg)d\mu(x)
				+\frac{\tau}{T}\int_{\{\rho_0>0\}}u\left(\frac{T}{\tau}\frac{\beta_{\tau}^{1/n}}{\rho_T(T(x))^{1/n}}\right)d\mu(x)
			\end{align}
			Now, we substitute back $u(\delta)=\delta^{n}U(\delta^{-n})$ and use the relation 
			\begin{equation}
				\frac{1}{\rho_T(T(x))}d\mu(x)=\frac{\rho_0(x)}{\rho_T(T(x))}d\mathscr{L}^n(x)=\det(\nabla T(x))d\mathscr{L}^n(x),
			\end{equation}
			to get
			\begin{align}
				&\frac{T-\tau}{T}\int_{\{\rho_0>0\}}u\left(\frac{T}{T-\tau}\frac{\beta_{T-\tau}^{1/n}}{\rho_0(x)^{1/n}}\right)d\mu(x)\\
				&\qquad+\frac{\tau}{T}\int_{\{\rho_0>0\}}u\left(\frac{T}{\tau}\frac{\beta_{\tau}^{1/n}}{\rho_T(T(x))^{1/n}}\right)d\mu(x)\\
				&=\left(\frac{T}{T-\tau}\right)^{n-1}\beta_{T-\tau}\int_{\{\rho_0>0\}}U\bigg(\frac{\rho_0(x)}{\beta_{T-\tau}}
				\left(\frac{T-\tau}{T}\right)^n\bigg)\frac{1}{\rho_0(x)}d\mu(x)\\
				&\qquad+\left(\frac{T}{\tau}\right)^{n-1}\beta_{\tau}\int_{\{\rho_0>0\}}U\bigg(\frac{\rho_T(T(x))}{\beta_{\tau}}
				\left(\frac{\tau}{T}\right)^n\bigg)\frac{1}{\rho_T(T(x))}d\mu(x)\\
				&=\left(\frac{T}{T-\tau}\right)^{n-1}\beta_{T-\tau}\int_{\mathbb{R}^n}U\bigg(\frac{\rho_0(x)}{\beta_{T-\tau}}
				\left(\frac{T-\tau}{T}\right)^n\bigg)d\mathscr{L}^n(x)\\
				&\qquad+\left(\frac{T}{\tau}\right)^{n-1}\beta_{\tau}\int_{\mathbb{R}^n}U\bigg(\frac{\rho_T(y)}{\beta_{\tau}}
				\left(\frac{\tau}{T}\right)^n\bigg)d\mathscr{L}^n(y).
			\end{align}
			This completes the proof of the entropic inequality.
		\end{proof}

\end{document}
